% ─────────────────────────────────────────────────────────────────
% CENG 467 Midterm – Question 3: Text Summarization
% Student ID: 310201050
% ─────────────────────────────────────────────────────────────────

\section{Question 3: Text Summarization}

\subsection{Dataset}

The CNN/DailyMail v3.0.0 corpus \cite{see2017get} is a large-scale
news-summarisation benchmark consisting of articles from CNN and the
Daily Mail paired with multi-sentence editor-written highlights. The
full corpus contains $\sim$287k training, 13k validation, and 11k test
articles. For this study, the first $200$ examples of the test split
are used (selected deterministically with seed 3102). The mean article
length in the subset is $555$ words; the mean reference summary length
is $34$ words across roughly 3 sentences. These characteristics motivate
setting LexRank's sentence budget to 3.

\subsection{Approaches}

\subsubsection{LexRank (Extractive, Graph-Based)}

LexRank \cite{erkan2004lexrank} represents each sentence in an article as
a node in an undirected graph, with edge weights equal to TF-IDF cosine
similarity between sentence pairs. A PageRank-style stationary
distribution is then computed, and the $k$ highest-ranking sentences are
returned verbatim as the summary. Because LexRank operates by selecting
existing sentences, it cannot introduce unfaithful content (every output
token comes from the source document), but it also cannot rephrase,
compress, or fuse information across sentences.

We use the implementation provided by the \texttt{sumy} library with the
English Snowball stemmer and standard English stopwords. Three sentences
are extracted per article, matching the average sentence count of the
gold-standard highlights.

\subsubsection{DistilBART (Abstractive, Encoder–Decoder Transformer)}

We use \texttt{sshleifer/distilbart-cnn-12-6}, a distilled BART
\cite{lewis2020bart} checkpoint with 12 encoder layers and 6 decoder
layers that has been pre-fine-tuned on CNN/DailyMail. The choice of
distilled weights is dictated by the M2 8\,GB memory budget. Generation
hyperparameters are summarised in Table~\ref{tab:bart_gen}.

\begin{table}[h]
\centering
\caption{DistilBART Generation Configuration}
\label{tab:bart_gen}
\begin{tabular}{ll}
\toprule
Parameter & Value \\
\midrule
Maximum input length        & 1024 sub-tokens (truncation) \\
Minimum / maximum output len & 30 / 130 sub-tokens \\
Decoding strategy           & Beam search, $B = 4$ \\
Length penalty              & 2.0 \\
No-repeat $n$-gram size     & 3 \\
Early stopping              & enabled \\
\bottomrule
\end{tabular}
\end{table}

\subsection{Evaluation Metrics}

Summaries are evaluated against the gold highlights using a multi-faceted
suite that captures complementary aspects of quality:

\begin{itemize}
  \item \textbf{ROUGE-1, ROUGE-2, ROUGE-L} \cite{lin2004rouge}: $F_1$ of
        unigram, bigram, and longest-common-subsequence overlap.
        Standard for summarisation but rewards lexical overlap rather
        than meaning.
  \item \textbf{BLEU} \cite{papineni2002bleu}: $n$-gram precision with
        brevity penalty (smoothing method 4). Borrowed from machine
        translation; tends to penalise abstractive paraphrasing.
  \item \textbf{METEOR} \cite{banerjee2005meteor}: incorporates stemming
        and synonym matching, partially correcting for ROUGE/BLEU's
        surface-level rigidity.
  \item \textbf{BERTScore F1} \cite{zhang2020bertscore}: token-level
        cosine similarity in contextual embedding space
        (\texttt{distilbert-base-uncased} backbone for memory
        efficiency). Captures semantic equivalence even under heavy
        paraphrasing.
\end{itemize}

\subsection{Results}

\begin{table}[h]
\centering
\caption{Summarisation Results on CNN/DailyMail (n=200 test articles)}
\label{tab:q3_results}
\begin{tabular}{lcc}
\toprule
Metric              & LexRank (extractive) & DistilBART (abstractive) \\
\midrule
ROUGE-1             & 0.2808 & \textbf{0.3538} \\
ROUGE-2             & 0.0939 & \textbf{0.1498} \\
ROUGE-L             & 0.1883 & \textbf{0.2580} \\
BLEU                & 0.0553 & \textbf{0.1043} \\
METEOR              & 0.3075 & \textbf{0.3370} \\
BERTScore F1        & 0.7776 & \textbf{0.8030} \\
\bottomrule
\end{tabular}
\end{table}

DistilBART outperforms LexRank on every single metric. The gap is
largest for $n$-gram-overlap metrics that reward lexical fidelity to the
reference (ROUGE-2 nearly $1.6\times$ higher; BLEU nearly $2\times$
higher). The gap on METEOR and BERTScore is narrower because both
metrics partially reward LexRank's verbatim copying of source-document
phrasings.

\subsection{Qualitative Analysis}

Three example articles are inspected in full in
\texttt{q3\_results.json} (indices 0, 50, 100). Three differences
between the extractive and abstractive outputs are consistently observed.

\paragraph{Fluency.}
LexRank summaries are by construction grammatical sentences from the
original article, so they are individually fluent. However, when the
three top-ranked sentences come from non-adjacent positions in the
source, the resulting concatenation reads as disjointed narrative ---
the test-set example at index 50 illustrates this: LexRank concatenates
three legal-discussion sentences about defamation that lose any
connection to the article's main narrative about the Rolling Stone
retraction. DistilBART, in contrast, produces a coherent paragraph
because the decoder generates it autoregressively under a smooth
language-model prior.

\paragraph{Information Coverage.}
DistilBART recovers the article's most newsworthy facts (\textit{``A new
trial was ordered in 2014 after firearms experts testified''} in example
100) that LexRank misses entirely, because such facts are spread across
multiple low-salience sentences in the source. LexRank's centrality-based
selection rewards repeated themes rather than rare informative content.

\paragraph{Factual Faithfulness.}
This is where the trade-off reverses. LexRank's outputs are extractively
faithful by definition. DistilBART, in contrast, occasionally exhibits
\emph{hallucination}: in example 50, the abstractive output attributes
the article to a named columnist (``Dean Obeidallah'') whose name does
not appear in the source preview. Such hallucinations are well
documented for abstractive encoder–decoder models trained with
maximum-likelihood objectives, and they remain a key open problem for
the field.

\subsection{Discussion}

ROUGE and BLEU favour systems that lexically overlap with the reference,
so LexRank --- which copies source sentences directly --- typically
scores respectably on these metrics despite generating less coherent
prose. However, the present results show that even on these surface
metrics, abstractive DistilBART wins decisively, because its
fine-tuning on CNN/DailyMail-style highlights teaches it to produce
output that closely matches the lexical distribution of human-written
summaries. DistilBART's lead widens further on METEOR and BERTScore,
both of which reward semantic alignment beyond surface form.

In terms of cost, LexRank requires only TF-IDF computation and a small
eigenvector calculation, runs entirely on CPU, and produces summaries
in milliseconds. DistilBART requires a fine-tuned $\sim$300\,MB model
and beam-search decoding that takes seconds per article. On readability
and coverage, DistilBART is clearly preferable. On faithfulness,
LexRank is strictly safer.

In production settings, the choice depends heavily on downstream
tolerance for hallucination: legal, medical, and financial
summarisation favour extractive (or strongly grounded abstractive)
approaches, whereas consumer-facing news digests can usually tolerate
the abstractive trade-off in exchange for fluency.